\subsection{布绕尔群的例子}

在下列三种情形，布绕尔群 \(\mathrm{Br}(\mathrm{K})\) 化归为单位元素：

\begin{enumerate}
\def\labelenumi{\alph{enumi})}
\tightlist
\item
  K 是可分闭的（VIII, 第 253 页，推论 1）。
\end{enumerate}

*b) K 是一个有限域（VIII, 第 357 页，推论 2）。

\begin{enumerate}
\def\labelenumi{\alph{enumi})}
\setcounter{enumi}{2}
\tightlist
\item
  K 具有性质 \((\mathrm{C}_1)\)（VIII, 第 357 页，注 2）。
\end{enumerate}

假设 K 是一个极大序域（VI, §2, No.~5, 第 25 页）。则 K 的布绕尔群是 2 阶循环群；其元素是 K 的类与 \((-1,0,-1)\) 型四元数 K-代数的类（III, §2, No.~5, 第 444 页与 VIII, 第 367 页，定理 1）。

假设 K 是一个局部紧的、非离散的交换拓扑域。若 K 不是连通的，则它是一个对离散赋值完备的域，且剩余域有限（Comm. Alg., VI, §9, No.~3, 第 433 页，定理 1），且存在一个从 Br(K) 到 Q/Z 的同构（VIII, 第 332 页，习题 17）。若 K 是连通的，则它同构于 R 或 C。域 R 的布绕尔群是 2 阶循环群。其非平凡元素是哈密顿四元数代数 H 的类（Gen.~Top., VIII, §1, No.~4, 第 104 页）；C 的布绕尔群的阶为 1。

